% !TEX root = ../thesis.tex
\chapter{Appendix}\label{ch:interpolation-lemma}\label{appendix:AM}
\section{Equivalence of continuous and discrete almost minimizing }
All (one–parameter) deformations are generated by vector fields compactly supported in the relevant
relative open set and tangent to $\partial M$ near $\partial M$. We always measure areas as
$\mass(\Gamma)=\mathcal{H}^n(\Gamma)$, and write $\partial_I\Omega$ for the interior boundary of a
finite–perimeter set $\Omega\subset M$.

\begin{definition}[Family of Surfaces]\label{def:family-of-surfaces}
A family $\{\Gamma_t\}_{t \in [0, 1]}$ of closed subsets of $M$ is a \textbf{generalized smooth family of surfaces} if:
\begin{enumerate}
    \item[(s1)] For each $t \in [0, 1]$, there exists a finite set $P_t \subset M$ such that $\Gamma_t$ is a smooth hypersurface with boundary in $M \setminus P_t$;

    \item[(s2)] For any compact set $K \subset M \setminus P_{t_0}$, we have $\Gamma_t \to \Gamma_{t_0}$ smoothly in $K$ as $t \to t_0$;

    \item[(s3)] (No concentration of mass) For every point $x \in M$:
    \begin{equation*}
        \limsup_{r \to 0}\sup_{t \in [0, 1]}\mathcal{H}^n(\Gamma_t\cap B_r(x)) = 0.
    \end{equation*}
\end{enumerate}
If, in addition to (s1)-(s3) above,
\begin{enumerate}
    \item [(s4)] The map $t \mapsto \mathcal{H}^n(\Gamma_t\cap M)$ is continuous, and $t \mapsto \Gamma_t$ is continuous in the Hausdorff topology,
\end{enumerate}
we moreover say that the generalized smooth family of surfaces $\{\Gamma_t\}_{t \in [0, 1]}$ is \textbf{continuous}.
\end{definition}

\begin{definition}[Sweepout (Technical Version for Interpolation)]\label{def:sweepout-appendix}
Let $U \subset M$ be a bounded relative open subset. A family $\{\Gamma_t\}_{t \in [0, 1]}$ is a \textbf{sweepout} of $U$ (in the sense of this appendix) if it satisfies (s1)-(s4) from Definition~\ref{def:family-of-surfaces} and there exists a family $\{\Omega_t\}_{t \in [0, 1]}$ of relative open sets in $M$ with finite measure such that:
\begin{enumerate}
    \item[(sw1)] $\Gamma_t \setminus \partial \Omega_t \subset P_t$ for all $t$, where $P_t$ is the finite singular set from (s1);

    \item[(sw2)] $\Omega_0 \cap U = \varnothing$ and $U \subset \Omega_1$;

    \item[(sw3)] $\mathcal{H}^{n + 1}(\Omega_t \triangle \Omega_s) \to 0$ as $|t - s| \to 0$.
\end{enumerate}
We call $\{\Omega_t\}$ the \textbf{corresponding family of open sets}.
The number $\mathcal{F}(\{\Gamma_t\}):=\sup_{t\in[0,1]}\mathcal{H}^n(\Gamma_t\cap M)$ is called the \textbf{height} of the sweepout.
\end{definition}

\begin{definition}[Continuous $\varepsilon$–a.m.\ in a relative open set]\label{def:cts-am-rel}\label{def:AM-rel}
Let $\varepsilon>0$ and $U\subset M$ be a relative open set. A relative boundary
$\Gamma=\partial_M\Omega\in\Zk$ is \emph{continuously $\varepsilon$–almost minimizing in $U$ with free boundary}
if there is no $C^0$ one–parameter family $\{\Gamma_t=\partial_M\Omega_t\}_{t\in[0,1]}\subset\Zk$
satisfying:
\begin{enumerate}[label=(\alph*)]
\item \emph{(sweepout regularity)} Properties \textup{(s1), (s2), (s3), (sw1), (sw3)} in Definition~\ref{def:sweepout-appendix};
\item \emph{(support in $U$)} $\Omega_t\setminus U=\Omega\setminus U$ for all $t$;
\item \emph{(uniform cost)} $\mass(\Gamma_t)\le \mass(\Gamma)+\varepsilon/8$ for all $t$;
\item \emph{(strict gain)} $\mass(\Gamma_1)\le \mass(\Gamma)-\varepsilon$.
\end{enumerate}
A sequence $\{\Gamma_k\}$ is (continuously) almost minimizing in $U$ if $\Gamma_k$ is
$\varepsilon_k$–a.m.\ with $\varepsilon_k\downarrow 0$.
\end{definition}

\begin{definition}[Discrete $(\varepsilon,\delta)$–a.m.\ in a relative open set]\label{def:dis-am-rel}
Let $\nu\in\{\mathcal F,\mathbf M,\mathbf F\}$ and $\varepsilon,\delta>0$.
Write $\mathscr A(U;\varepsilon,\delta;\nu)$ for the set of all
$\Gamma=\partial_M\Omega\in\Zk$ with the following property:

for every finite sequence $\Gamma^0,\Gamma^1,\dots,\Gamma^m\in\Zk$ with
\begin{enumerate}[label=(\roman*)]
\item \emph{(support in $U$)} $\Omega^j\setminus U=\Omega\setminus U$ for $j=0,\dots,m$;
\item \emph{(small steps)} $\nu(\Gamma^{j+1}-\Gamma^{j})\le \delta$ for $j=0,\dots,m-1$;
\item \emph{(no overshoot)} $\mass(\Gamma^{j})\le \mass(\Gamma)+\delta$ for $j=1,\dots,m$,
\end{enumerate}
one necessarily has the terminal lower bound
\[
\mass(\Gamma^{m})\ \ge\ \mass(\Gamma)-\varepsilon.
\]
We say that a varifold $V$ is \emph{discrete a.m.\ in $U$} if there are
$\varepsilon_i\downarrow 0$, $\delta_i\downarrow 0$ and $\Gamma_i\in\mathscr A(U;\varepsilon_i,\delta_i;\mathcal F)$
such that $\mathbf F(|\Gamma_i|,V)<\varepsilon_i$.
\end{definition}

\begin{lemma}[Relative interpolation]\label{lem:relative-interpolation}
Fix $L>0$. For every $\eta>0$ there exists $\delta>0$ with the following property.
If $\Omega_0,\Omega_1\subset M$ are finite–perimeter sets with $\Omega_0\subset\Omega_1$,
$\Per(\Omega_i)\le L$ and $\mathcal{H}^{n+1}\big((\Omega_1\setminus\Omega_0)\cap U\big)\le \delta$,
then there exists a nested, $\mathcal F$–continuous family
$\{\partial_M\Omega_t\}_{t\in[0,1]}$ supported in $U$ with
\[
\Per(\Omega_t)\ \le\ \max\{\Per(\Omega_0),\Per(\Omega_1)\}+\eta\qquad\text{for all }t.
\]
All flows used in the construction are tangent to $\partial M$ near $\partial M$.
\end{lemma}

\begin{proof}[Idea]
Cover $(\Omega_1\!\setminus\!\Omega_0)\cap U$ by finitely many small balls whose boundary spheres
meet $\partial_M\Omega_i$ transversely and have small relative area (coarea).
Apply the preliminary modification \ref{lem:preliminary-modification} to add/remove these pieces
one by one inside the balls, and use Lemma~\ref{lem:existence-nested-between-sets} to produce the
nested homotopies with the stated perimeter bound.
\end{proof}

\begin{remark}
Lemma~\ref{lem:relative-interpolation} is the relative version of the interpolation lemma in \cite[Lemma~B.1]{LZ17}, adapted to the free boundary setting.
\end{remark}

\begin{proposition}[Continuous $\Rightarrow$ discrete, quantitative]\label{prop:cts-to-dis}
There exists an absolute constant $c_0\in(0,1)$ with the following property.
If $\Gamma=\partial_M\Omega$ is continuously $\varepsilon$–a.m.\ in $U$,
then $\Gamma\in\mathscr A(U;\varepsilon,\delta;\mathcal F)$ for every
$\delta\le c_0\,\varepsilon$.
\end{proposition}

\begin{proof}
Let $\{\Gamma^j\}_{j=0}^m$ be any discrete competitor as in Definition~\ref{def:dis-am-rel}
with parameter $\delta>0$ to be fixed.
For each pair $(\Gamma^j,\Gamma^{j+1})$ we apply the relative local modification
(Preliminary Modification, §6.2) to decompose their difference inside $U$ into a finite sequence of
nested elementary additions/removals supported on small balls, with the property that every
intermediate slice has area bounded by
\[
\max\{\mass(\Gamma^j),\mass(\Gamma^{j+1})\}+C\,\delta,
\]
for a dimensional $C=C(n)$.
Using Lemma~\ref{lem:relative-interpolation} on each elementary step and concatenating in time,
we produce a $C^0$ family $\{\widehat\Gamma_t\}_{t\in[0,1]}$ connecting $\Gamma^0=\Gamma$ to
$\Gamma^m$, supported in $U$, and satisfying
\[
\sup_t\mass(\widehat\Gamma_t)\ \le\ \max_{0\le j\le m}\mass(\Gamma^j)+C\,\delta
\ \le\ \mass(\Gamma)+\delta+C\,\delta.
\]
Choosing $\delta\le c_0\varepsilon$ with $c_0>0$ sufficiently small (e.g.\ $c_0\le 1/16C$),
we ensure
\[
\sup_t\mass(\widehat\Gamma_t)\ \le\ \mass(\Gamma)+\varepsilon/8.
\]
Hence $\{\widehat\Gamma_t\}$ is an admissible competitor in the sense of
Definition~\ref{def:cts-am-rel}. By the continuous $\varepsilon$–a.m.\ property we must have
\(\mass(\Gamma^m)\ge \mass(\Gamma)-\varepsilon\).
This proves $\Gamma\in \mathscr{A}(U;\varepsilon,\delta;\mathcal F)$ for all $\delta\le c_0\varepsilon$.
\end{proof}

\begin{corollary}[Discrete a.m.\ for the min–max limit]\label{cor:dis-am-limit}
Let $\{\Gamma_k\}$ be a min–max sequence which is continuously $\varepsilon_k$–a.m.\ in $U$
with $\varepsilon_k\downarrow 0$, and suppose $|\Gamma_k|\to V$ as varifolds.
Then $V$ is discrete almost minimizing in $U$ with free boundary.
\end{corollary}

\begin{proof}
By Proposition~\ref{prop:cts-to-dis}, for each $k$ we have
$\Gamma_k\in \mathscr{A}(U;\varepsilon_k,\delta_k;\mathcal F)$ with
$\delta_k=c_0\varepsilon_k\downarrow 0$. Since $\mathbf F(|\Gamma_k|,V)\to 0$,
the definition of discrete a.m.\ for varifolds is satisfied.
\end{proof}

\begin{remark}[Pairs of annuli]
All statements above hold verbatim when $U$ is replaced by a ball, an annulus, or a \emph{pair}
of well–separated annuli $(A_1,A_2)$ (in the sense of $\mathrm{CO}(\mathcal A)$ used in the text):
the deformations are supported in $A_1\cup A_2$, and the proofs are unchanged.
\end{remark}

\section{Fermi Coordinates Near the Boundary}\label{sec:fermi-coordinates}

In this section, we provide a self-contained exposition of \emph{Fermi coordinates} near a boundary
point of a Riemannian manifold with boundary. These coordinates are particularly well-suited for
studying free-boundary minimal hypersurfaces because they naturally encode the orthogonality
condition at $\partial M$. The presentation follows the framework developed by Li--Zhou~\cite{LZ17}
and is adapted to our setting.

\subsection{Definition and Basic Properties}

\begin{definition}[Fermi coordinates]\label{def:fermi-coordinates}
Let $p \in \partial M$. Suppose $(x_1,x_2,\cdots,x_n)$ is a geodesic normal coordinate system for
$\partial M$ centered at $p$, defined in a small neighborhood of $p$ in $\partial M$. Let
$t=\dist_M(\cdot,\partial M)$ denote the signed distance function from $\partial M$, which is
smooth in a sufficiently small relatively open neighborhood of $p$ in $M$. Then
$(x_1,x_2,\cdots,x_n,t)$ is called a \emph{Fermi coordinate system} for $(M,\partial M)$ centered
at $p \in \partial M$.

We define the \emph{Fermi distance function from $p$} on a relatively open neighborhood of $p$ in
$M$ by
\begin{equation}\label{eq:fermi-distance}
\tilde r=\tilde r_p(q):=\sqrt{x_1^2+\cdots+x_n^2+t^2}.
\end{equation}
The \emph{Fermi exponential map at $p$}, denoted $\tilde{\exp}_p$, gives the Fermi coordinate
system and is defined on a half-ball in $\R^{n+1}_+ \cong T_pM$ centered at the origin. Explicitly,
$\tilde{\exp}_p(x_1,\cdots,x_n,t)$ is the point in $M$ whose Fermi coordinates are
$(x_1,\cdots,x_n,t)$.
\end{definition}

The key feature of Fermi coordinates is that the $t$-direction is always orthogonal to $\partial M$,
which makes them ideal for analyzing free-boundary conditions.

\begin{lemma}[Metric and connection in Fermi coordinates]\label{lem:fermi-metric}
Let $g_{ab}$ and $\Gamma_{ab}^c$ be the components of the metric and Christoffel symbols in Fermi
coordinates $(x,t)=(x_1,\cdots,x_n,t)$ centered at $p \in \partial M$. Then for $i,j=1,\cdots,n$,
\begin{equation}\label{eq:fermi-metric-properties}
g_{tt} \equiv 1, \quad g_{it} \equiv 0, \quad \Gamma_{it}^t=\Gamma_{tt}^i=\Gamma_{tt}^t \equiv 0.
\end{equation}
Moreover, there exist constants $r_{\mathrm{Fermi}},C>0$ depending only on the geometry of $M$
(specifically, bounds on the curvature and second fundamental form of $\partial M$) such that for
$\tilde r \in [0,r_{\mathrm{Fermi}})$,
\begin{equation}\label{eq:fermi-estimates}
\begin{split}
|g_{ij}(x,t)-\delta_{ij}| &\leq C \tilde r, \\
|\Gamma_{ij}^k(x,t)| &\leq C \tilde r, \\
|\Gamma_{it}^j(x,t)|+|\Gamma_{ij}^t(x,t)| &\leq C.
\end{split}
\end{equation}
\end{lemma}

\begin{proof}
Since $t=\dist_M(\cdot,\partial M)$ is the distance function from the boundary, its gradient has
unit length and is orthogonal to level sets. This immediately gives $g_{tt}\equiv 1$ and
$g_{it}\equiv 0$. The vanishing of the Christoffel symbols $\Gamma_{it}^t$, $\Gamma_{tt}^i$, and
$\Gamma_{tt}^t$ follows from the fact that $\frac{\partial}{\partial t}$ has constant length and
is orthogonal to $\frac{\partial}{\partial x_i}$.

For the quantitative estimates, observe that $\frac{\partial}{\partial x_i}$ are Jacobi fields
along the geodesics $\gamma_x(t)=(x,t)$ emanating orthogonally from $\partial M$. They satisfy
the Jacobi equation
\[
\nabla_{\frac{\partial}{\partial t}} \nabla_{\frac{\partial}{\partial t}} \frac{\partial}{\partial x_i}
= -R\left(\frac{\partial}{\partial t},\frac{\partial}{\partial x_i}\right) \frac{\partial}{\partial t},
\]
where $R$ is the Riemann curvature tensor of $M$.

By standard estimates for geodesic normal coordinates on $\partial M$, for $|x|<r_{\mathrm{Fermi}}$
we have
\begin{equation}\label{eq:g-ij-boundary}
|g_{ij}(x,0)-\delta_{ij}| \leq C |x|,
\end{equation}
as $\partial M$ has bounded sectional curvature. Along the $t$-direction, we have
$\partial_t g_{ij}(x,0)=h_{ij}(x)$, where $h_{ij}$ are the components of the second fundamental
form of $\partial M$ in the coordinates $(x_1,\cdots,x_n)$. Using the Jacobi equation, one derives
the Riccati-type equation (see~\cite[Lemma~2.2]{Marques05})
\[
\partial_t^2 g_{ij}=-2 R_{titj} +2g^{rs} (\partial_t g_{ir}) (\partial_t g_{sj}).
\]
Since the curvature is bounded, integrating this equation yields for $t < r_{\mathrm{Fermi}}$,
\begin{equation}\label{eq:g-ij-evolution}
|g_{ij}(x,t) - g_{ij}(x,0)| \leq C t.
\end{equation}
Combining \eqref{eq:g-ij-boundary} and \eqref{eq:g-ij-evolution} gives
$|g_{ij}(x,t)-\delta_{ij}|\leq C(|x|+t)\leq C\tilde r$. The estimates for the Christoffel symbols
follow from similar arguments using bounds on the derivatives of the curvature tensor.
\end{proof}

\subsection{Geometric Properties of the Fermi Distance Function}

\begin{lemma}[Gradient and Hessian of the Fermi distance]\label{lem:fermi-distance-estimates}
The Fermi distance function $\tilde r$ is smooth (where it is defined) and there exist constants
$r_{\mathrm{Fermi}},C>0$ depending only on the geometry of $M$ such that for
$\tilde r \in [0,r_{\mathrm{Fermi}})$,
\begin{enumerate}[label=(\roman*)]
\item $\left\|\nabla \tilde r - \frac{\partial}{\partial \tilde r}\right\|_g \leq C \tilde r$,
\item $\left\| \Hess \tilde r - \frac{1}{\tilde r} g_{\tilde r} \right\|_g \leq C$,
\end{enumerate}
where $g_{\tilde r}$ denotes the induced metric on the sphere of radius $\tilde r$ in Euclidean
space (i.e., the round metric scaled by $\tilde r$).
\end{lemma}

\begin{proof}
The gradient of $\tilde r$ in Fermi coordinates is given by
\begin{equation}\label{eq:grad-fermi-distance}
\nabla \tilde r = \sum_{i,j=1}^n g^{ij}(x,t) \frac{x_i}{\tilde r} \frac{\partial}{\partial x_j}
+\frac{t}{\tilde r} \frac{\partial}{\partial t}.
\end{equation}
Using $g^{ij}=\delta_{ij}+O(\tilde r)$ from Lemma~\ref{lem:fermi-metric}, we obtain
\[
\nabla \tilde r = \frac{1}{\tilde r}\left(\sum_{i=1}^n x_i \frac{\partial}{\partial x_i}
+ t \frac{\partial}{\partial t}\right) + O(\tilde r)
= \frac{\partial}{\partial \tilde r} + O(\tilde r),
\]
proving (i).

For (ii), direct computation in Fermi coordinates yields
\[
\Hess \tilde r \left(\frac{\partial}{\partial t}, \frac{\partial}{\partial t}\right)
= \frac{\tilde r^2-t^2}{\tilde r^3}, \quad
\left| \Hess \tilde r \left(\frac{\partial}{\partial x_i}, \frac{\partial}{\partial t}\right)
+ \frac{x_i t}{\tilde r^3} \right| \leq C,
\]
\[
\left| \Hess \tilde r \left(\frac{\partial}{\partial x_i}, \frac{\partial}{\partial x_j}\right)
- \frac{\delta_{ij} \tilde r^2-x_i x_j}{\tilde r^3}\right| \leq C.
\]
These formulas show that $\Hess \tilde r$ agrees with the Hessian of the Euclidean distance
function up to bounded error, which is precisely the statement (ii).
\end{proof}

\subsection{Fermi Balls and Convexity}

\begin{definition}[Fermi balls and spheres]\label{def:fermi-balls}
For each $p \in \partial M$, we define the \emph{Fermi half-ball} and \emph{Fermi half-sphere}
of radius $r$ centered at $p$ respectively by
\[
\tilde B^+_r(p):= \{ q \in M : \tilde r_p(q) < r \}, \qquad
\tilde S^+_r(p):=\{ q \in M : \tilde r_p(q) = r \}.
\]
\end{definition}

The following lemma summarizes the key geometric properties of Fermi half-balls that are most
relevant for our applications.

\begin{lemma}[Convexity and comparison with Euclidean balls]\label{lem:fermi-balls-properties}
There exists a constant $r_{\mathrm{Fermi}}>0$, depending only on the geometry of $M$, such that
for all $0<r<r_{\mathrm{Fermi}}$ and $p\in\partial M$,
\begin{enumerate}[label=(\alph*)]
\item $\tilde S^+_r(p)$ is a smooth hypersurface meeting $\partial M$ orthogonally,
\item $\tilde B^+_r(p)$ is a relatively convex domain in $M$,
\item $B_{r/2}(p) \cap M \subset \tilde B^+_r(p) \subset B_{2r}(p) \cap M$,
\end{enumerate}
where $B_s(p)$ denotes the Euclidean ball of radius $s$ centered at $p$ in the ambient space
$\mathbb{R}^L$.
\end{lemma}

\begin{proof}
Properties (a) and (c) follow directly from the estimates in Lemmas~\ref{lem:fermi-metric}
and~\ref{lem:fermi-distance-estimates}. For (b), the relative convexity of $\tilde B^+_r(p)$
follows from the estimate (ii) in Lemma~\ref{lem:fermi-distance-estimates}: the Hessian of
$\tilde r$ is positive definite (up to a bounded perturbation from the Euclidean case), which
implies that sublevel sets of $\tilde r$ are convex for sufficiently small $r$.
\end{proof}

\begin{remark}[Applications to free-boundary theory]
Fermi coordinates are particularly useful in the regularity theory for free-boundary minimal
hypersurfaces. The orthogonality condition $\Gamma\perp\partial M$ translates to the statement
that $\Gamma$ is a graph over a domain in $\{t=0\}$ near $\partial M$, and the free-boundary
condition becomes a Neumann boundary condition for the graphing function. This perspective is
exploited extensively in the work of Li--Zhou~\cite{LZ17} and in our application of their
regularity results in Chapter~\ref{ch:convergence-of-min-max}.
\end{remark}

\section{Wickramasekera's Regularity Theory for Stable Varifolds}\label{sec:wickramasekera-detailed}

In this section, we provide a detailed exposition of Wickramasekera's regularity theory for stable
codimension 1 integral varifolds~\cite{Wic14}, which plays a crucial role in establishing the
smoothness of minimal hypersurfaces arising from non-excessive min-max procedures in
Chapter~\ref{ch:alternative-regularity}.

\subsection{Stationary and Stable Varifolds}

We begin by recalling the key definitions. Let $V$ be an integral $n$-varifold on a smooth
$(n+1)$-dimensional Riemannian manifold $N$.

\begin{definition}[Stationary varifold]\label{def:stationary-varifold}
$V$ is \emph{stationary} if it has zero first variation with respect to area, i.e., for any
compactly supported $C^1$ vector field $\psi$ on $N$,
\[
\int_{N \times G_n} \mathrm{div}_{S} \,\psi(X) \, dV(X, S) = 0,
\]
where $G_n$ denotes the space of unoriented $n$-planes in the tangent bundle of $N$, and
$\mathrm{div}_S$ denotes the tangential divergence on the plane $S$.
\end{definition}

For a stationary varifold $V$, the \emph{regular part} $\mathrm{reg}\, V$ is the set of points where
the support of $V$ is a smooth embedded hypersurface, and the \emph{singular set}
$\mathrm{sing}\, V$ is its complement in $\mathrm{spt}\,\|V\|$. By Allard's regularity
theorem~\cite{AW}, $\mathrm{reg}\, V$ is always a dense, open subset of $\mathrm{spt}\,\|V\|$ when
$V$ is stationary.

\begin{definition}[Stable varifold]\label{def:stable-varifold}
A stationary integral $n$-varifold $V$ is \emph{stable} if its regular part satisfies the
\emph{stability inequality}: for every open set $\Omega \subset N$ with sufficiently small
singular set (e.g., $\mathrm{sing}\, V \cap \Omega = \emptyset$ when $n \leq 6$, or
${\mathcal H}^{n-7+\gamma}(\mathrm{sing}\, V \cap \Omega) = 0$ for all $\gamma > 0$ when $n \geq 7$),
\begin{equation}\label{eq:stability-detailed}
\int_{\mathrm{reg} \, V \cap \Omega} |A|^{2}\zeta^{2} \,d{\mathcal H}^{n} \leq
\int_{\mathrm{reg} \, V \cap \Omega} |\nabla \zeta|^{2} \, d{\mathcal H}^{n}
\quad \forall \; \zeta \in C^{1}_{c}(\mathrm{reg} \, V \cap \Omega),
\end{equation}
where $A$ denotes the second fundamental form of $\mathrm{reg}\, V$, $|A|$ its norm, and $\nabla$
the gradient on $\mathrm{reg}\, V$.
\end{definition}

Equivalently, stability means that $V$ has non-negative second variation of area under normal
deformations supported away from the singular set. The stability inequality~\eqref{eq:stability-detailed}
is a quantitative expression of this variational property.

\subsection{The $\alpha$-Structural Hypothesis}

The key geometric hypothesis in Wickramasekera's theorem is the $\alpha$-Structural Hypothesis (Definition~\ref{defn:alpha-structural} in Chapter~\ref{ch:alternative-regularity}), which we recall here for convenience:

\begin{quote}
For $\alpha \in (0, 1/2)$, a stationary integral varifold $V$ satisfies the \emph{$\alpha$-Structural Hypothesis} if no singular point of $V$ has a neighborhood in which $V$ corresponds to a union of embedded $C^{1,\alpha}$ hypersurfaces-with-boundary meeting only along their common $C^{1,\alpha}$ boundary, with multiplicity a constant positive integer on each constituent hypersurface.
\end{quote}

In other words, the $\alpha$-Structural Hypothesis rules out the possibility that the varifold
locally looks like three or more sheets meeting along a common $(n-1)$-dimensional junction.

\begin{remark}
The $\alpha$-Structural Hypothesis is implied by either of the following conditions:
\begin{enumerate}[label=(\alph*)]
\item $V$ corresponds to an absolutely area-minimizing rectifiable current with no boundary.
\item The singular set has vanishing $(n-1)$-dimensional Hausdorff measure:
${\mathcal H}^{n-1}(\mathrm{sing}\, V) = 0$.
\end{enumerate}
Condition (b) is the weakest size hypothesis on the singular set that guarantees the regularity
conclusions below. No hypothesis in terms of Hausdorff measure of dimension $> n-1$ can suffice,
as shown by examples of transversely intersecting hyperplanes.
\end{remark}

\subsection{Wickramasekera's Main Theorem}

We can now state the complete version of Wickramasekera's regularity and compactness theorem:

\begin{theorem}[Wickramasekera's Regularity and Compactness Theorem~\cite{Wic14}]
\label{thm:wickramasekera-full}
Let $V$ be a stable integral $n$-varifold on a smooth $(n+1)$-dimensional Riemannian manifold $N$.
Suppose $V$ satisfies the $\alpha$-Structural Hypothesis for some $\alpha \in (0,1/2)$. Then:
\begin{enumerate}[label=(\roman*)]
\item \textbf{Regularity:}
\begin{itemize}
\item If $2 \leq n \leq 6$: $\mathrm{sing}\, V = \emptyset$ (complete regularity).
\item If $n = 7$: $\mathrm{sing}\, V$ is discrete.
\item If $n \geq 8$: ${\mathcal H}^{n-7+\gamma}(\mathrm{sing}\, V) = 0$ for all $\gamma > 0$.
\end{itemize}
In all cases, $V$ corresponds to a smooth embedded hypersurface away from the singular set, with
multiplicity a constant positive integer on each connected component.

\item \textbf{Compactness:} For any fixed $\alpha \in (0,1/2)$, each mass-bounded subset of the
class of stable codimension 1 integral varifolds satisfying the $\alpha$-Structural Hypothesis
is compact in the varifold convergence topology.
\end{enumerate}
\end{theorem}

The dimension estimates on the singular set in part (i) are optimal, as demonstrated by well-known
examples of stable minimal cones (e.g., the cone over
$\mathbf{S}^3(1/\sqrt{2}) \times \mathbf{S}^3(1/\sqrt{2}) \subset \mathbf{R}^8$ has an isolated
singularity at the origin).

\subsection{Application to Non-Excessive Min-Max Theory}

The relevance of Wickramasekera's theorem to our work in Chapter~\ref{ch:alternative-regularity} is as
follows. The Chodosh--Liokumovich--Spolaor theory~\cite{CLS22} produces a stationary varifold from
a non-excessive sweepout, but does not automatically guarantee integrality. Our contribution is to
show that in the multiplicity-one case, no mass cancellation occurs in the limit, which implies:
\begin{enumerate}[label=(\alph*)]
\item The limiting varifold is integral and rectifiable.
\item The singular set has vanishing $(n-1)$-dimensional Hausdorff measure.
\item The $\alpha$-Structural Hypothesis is satisfied.
\end{enumerate}
Given these properties, together with the stability implied by the non-excessive condition,
Theorem~\ref{thm:wickramasekera-full} immediately yields smoothness of the resulting minimal
hypersurface when $n \leq 6$.

\subsection{Key Ideas in the Proof}

While we do not reproduce the full proof of Theorem~\ref{thm:wickramasekera-full} (which occupies
a substantial monograph), we outline the main ideas:

\begin{enumerate}
\item \textbf{Sheeting Theorem:} Wickramasekera proves that if a stable varifold satisfying the
$\alpha$-Structural Hypothesis is weakly close to a hyperplane of multiplicity $q$, it must locally
decompose into $q$ disjoint smooth graphs (``sheets'') over the hyperplane.

\item \textbf{Minimum Distance Theorem:} No such varifold can be close to a cone consisting of
three or more half-hyperplanes meeting along a common axis.

\item \textbf{Inductive argument:} The proof proceeds by induction on the multiplicity, using blow-up
analysis and refined excess-decay estimates to control the geometry near potential singularities.

\item \textbf{Use of stability:} The stability inequality~\eqref{eq:stability-detailed} is used
crucially throughout, particularly in regions where the Hausdorff dimension of the singular set
is sufficiently small.
\end{enumerate}

A key technical point is that Wickramasekera's proof does \emph{not} assume a priori that the
singular set has finite $(n-2)$-dimensional Hausdorff measure (as in the earlier work of
Schoen--Simon~\cite{SS}). The $\alpha$-Structural Hypothesis is a purely geometric condition that
does not directly imply any bounds on the size of the singular set, yet it suffices to conclude
optimal regularity.

\begin{remark}[Comparison with area-minimizing theory]
For absolutely area-minimizing codimension 1 rectifiable currents, the same regularity conclusions
were established through the combined work of De Giorgi~\cite{DG}, Reifenberg, Fleming, Almgren,
Simons~\cite{SJ}, and Federer. Wickramasekera's theorem subsumes this classical theory: if $T$ is
an area-minimizing current with no boundary, it automatically satisfies the $\alpha$-Structural
Hypothesis (since no union of three or more sheets can be area-minimizing).
\end{remark}
